\chapter{Den fyrste reforma}

\begin{motto}
Dei sette seg føre å gje forma eit fundament.\\ Dei kom heim med ein katalog over ornament og ein dom over smak.
\end{motto}

I 1851 reiste det seg ein bygning av glas og jarn i Hyde Park, og bak glaset stilte verda ut seg sjølv. Den store utstillinga samla under eitt tak alt den industrielle forma hadde produsert: maskinar, møblar, tekstilar, prydgjenstandar frå heile imperiet og frå nasjonane det måle seg mot. Ein krins av britiske reformatorar gjekk gjennom hallane, såg på det utstilte, og fann det avskyeleg. Det dei såg, var den industrielle forma sleppt laus utan ein standard: maskina kunne no støype og prege ornament i ubunden mengd, og fabrikantane hadde gjort nett det, dekt kvar flate med pynt, etterlikna eitt materiale i eit anna, sett rokokko på dampmaskinar og naturalistiske blomar på teppe ein gjekk på. Jeffrey Auerbach har synt kor mangetydig sjølve hendinga var, på same tid nasjonal sjølvfeiring, imperial mønstring og estetisk sjokk \autocite{auerbach1999}. For oss tel det siste. For fyrste gong stod nokon andsynes den industrielle forma i full breidd og spurde: kva er det vi har laga, og kvifor er det så stygt.

Spørsmålet var rett. Det reformatorane ikkje såg, var kor stort det var. Dei las styggleiken som eit moralsk samanbrot, og dei hadde ikkje heilt urett, for noko hadde brote saman: den innebygde formkvaliteten handverket heldt oppe så lenge tenking og laging var eitt. Men dei sjølve stod i tomrommet etter det same skiljet, og oppgåva dei tok på seg, å gje den industrielle forma ein grunn, var av eit anna slag enn ho dei trudde dei tok. Dei trudde dei skulle reformere ein smak. Dei stod framfor det å byggje eit fundament. Dei gjorde det fyrste og kalla det det andre, og resten av kapitlet er den same hendinga om att: ein autoritet som får fysisk form, og som kollapsar i fyrste møte med motstand, fordi det aldri fanst nokon grunn under han utanom den som bar han.

\section{Dommen utan grunn: Cole og skammekammeret}

Mannen i sentrum var Henry Cole, embetsmann med ein sjeldan kombinasjon av estetisk overtyding og administrativ kraft. Cole hadde vore med å organisere utstillinga, og i åra kring henne bygde han eit heilt apparat for å betre den britiske smaken: tidsskriftet \textit{Journal of Design and Manufactures}, ein sentralskule, og etter kvart samlingane som skulle bli til South Kensington. Penny Sparke reknar dette mellom dei fyrste samordna statlege forsøka i historia på å heve den industrielle formkulturen \autocite{sparke2004}. Apparatet var verkeleg. Spørsmålet er kva slag kunnskap det hadde å dele ut, og det vart prøvd offentleg, ein einaste gong, med eit utfall Cole ikkje hadde rekna med.

I 1852 opna Cole sitt nye Department of Practical Art ei undervisningssamling ved Marlborough House, og i henne ein avdeling med tittelen «Examples of False Principles in Decoration». Her stilte Richard Redgrave, Cole sin næraste mann, ut innkjøpte samtidsvarer reformatorane meinte braut mot god smak: teppe med naturalistiske blomar ein gjekk på, gasslampar forma som blomar, tapet med perspektivisk djupn på ei flate som var flat. Kvar gjenstand bar ein dom og ei grunngjeving, kva «falskt prinsipp» han braut. Folket døypte avdelinga straks om til «the Chamber of Horrors» og straumde til. Her, om nokon stad, skulle reforma demonstrere kunnskapen sin offentleg: dette er stygt, og her er grunnen.

To ting hende. Publikum oppdaga at fleire av skrekkdøma var ting dei sjølve åtte og likte, og lét seg fortelje av embetsmenn at smaken deira var falsk. Og somme produsentar vart så rasande over å sjå varene sine hengde ut ved namn at det vart eit ope ordskifte, og avdelinga måtte byggjast om og dempast ned alt etter kort tid. Charles Dickens fanga episoden i ei satire i \textit{Household Words}, «A House Full of Horrors», om ein stakkar ved namn Crumpet som hadde lært reglane så grundig at han ikkje lenger kunne nyte nokon ting: kvar tapetrose vart eit brot på eit prinsipp, kvar teppeblome ein forbrytelse. Då Crumpet spør kvifor tapetrosa er falsk, finst det ikkje noko svar som ikkje til sjuande og sist er «fordi dei lærde fann henne falsk». Adrian Forty har lese heile skammekammeret som ei øving i klassemakt forkledd som estetikk, dei rette folka som lærte dei andre kva dei skulle mislike \autocite{forty1986}.

Episoden er liten, men ho er ei måling, og målinga er presis. Reforma sin kunnskap rakk akkurat så langt som autoriteten til Cole-krinsen, og ikkje ein tomme lenger. Då autoriteten møtte motstand frå folk som likte sine eigne ting og produsentar som selde dei, hadde reforma ingenting å falle tilbake på utan meir autoritet. Eit fundament hadde tolt protesten, fordi det hadde hatt ein grunn som stod uavhengig av kven som likte kva. Smaken tolte han ikkje. Skammekammeret måtte falle saman fordi det aldri var anna enn ein autoritet som hadde fått fysisk form. Dette er det boka kallar \emph{sviktsignaturen}: ein strid som ikkje kan avgjerast utan rang, frist eller stemme, fordi det manglar ein grunn over begge partar. Ein støyter på han her fyrste gong, og ein kjem til å støyte på han kvar gong reforma blir sett på prøve.

Ein ting til er verdt å feste, fordi det knyter kapitlet til skiljet boka opna med. Då reformatorane kalla ei vare «falsk», dømde dei overflata: det etterlikna materialet, det påklistra ornamentet, parafinlampen som gav seg ut for marmor. Men dei spurde aldri kvifor maskina øste ut nett dette i slik mengd. Svaret låg under dommen, i skiljet mellom form og følgje: den som teikna bar ikkje konsekvensen, og den som kjøpte kunne ikkje sjå skilnaden før heime. Den overlessa vara var ikkje synd hjå fabrikanten, ho var den rasjonelle forma i ein marknad som løna det å sjå dyr ut og kosta lite. Dommen var rett i det han mislikte og blind for kvifor det fanst. Han nådde overflata av vara og aldri ned til skiljet som produserte henne.

\section{Pugin: ein moral kledd som ein lov}

Bak smaksreforma låg ein djupare figur, og han er det næraste 1800-talet kom eit prinsipp om form. Augustus Welby Northmore Pugin var arkitekt, konvertitt og polemikar, og han gjorde noko ingen smaksdommar gjorde: han kopla form til moral og hevda at det fanst sanne prinsipp for korleis ting skulle lagast. I \textcite{pugin1836contrasts} stilte han mellomalderens bygningar opp mot samtidas og las forfallet i forma som eit forfall i samfunnet og trua. Og i \textcite{pugin1841true} formulerte han dei to setningane heile den komande designhistoria skulle leve på: at det ikkje skal finnast trekk ved ein bygning som ikkje er naudsynt for tilhøve, konstruksjon eller sømd, og at all ornamentering skal vere ein utsmykking av den essensielle konstruksjonen. Endeleg ser det ut som eit fundament. Ikkje berre ein dom om at noko er stygt, men eit prinsipp ein kan utleie det rette frå.

Granskar ein prinsippet, viser det seg å vere ein moral kledd som ein lov. Pugin sin grunn for at forma skulle vere slik, var ikkje ein teori om form, han var ein teologi. Det rette var det kristne, det gotiske, det mellomalderlege, fordi desse var sanne i religiøs meining, ikkje fordi noka analyse av form synte at dei var betre. «Ærlegdom» i konstruksjonen var ein dygd, ikkje ein dokumentert eigenskap; å skjule eit bjelkelag var ei løgn i moralsk forstand, ikkje ein feil i nokon målbar. Eit slikt prinsipp kan ikkje etterprøvast; det kan berre delast eller avvisast. Skift trua, og prinsippet fell. Det er nett dette som hender når seinare modernistar arvar ærlegdomsmoralen utan teologien og kallar han funksjon: dei held att forma, men har ingen grunn under han som ikkje sjølv er ein smak. Adolf Loos skulle drive same draget til sin ytste kant, då han eit halvt hundreår seinare kopla ornamentet til primitivitet og dekadanse og gjorde fråvêret av pynt til ein moralsk dom med motsett forteikn av Pugin sin \autocite{loos1908ornament}. Pugin lova eit prinsipp og leverte ein moral, og arven hans var stor nett fordi ein moral lèt seg flytte fritt frå éin grunn til ein annan.

\section{Jones: katalogen som heitte grammatikk}

Det mest ambisiøse forsøket bar tittelen eit fundament ville fortent. I 1856 gav Owen Jones ut \textit{The Grammar of Ornament}, eit praktverk som samla ornament frå alle kulturar og epokar, egyptisk, gresk, maurisk, keltisk, i fargelagde plansjar av ei prakt som framleis verkar \autocite{jones1856grammar}. Han ramma verket inn med trettisju proposisjonar, allmenne reglar for korleis ornament skulle byggjast: at konstruksjonen skal dekorerast og dekorasjonen aldri konstruerast, at all ornament skal byggje på geometri, at blomar skal stiliserast og ikkje etterliknast. Og han kalla det ein grammatikk. Han lova, i sjølve tittelen, det faget framleis manglar: eit sett reglar som genererer rett form, slik ein grammatikk genererer rette setningar.

Men ein grammatikk er generativ og forklarande. Ho fortel deg kvifor ei velforma setning er velforma, og let deg lage uendeleg mange du aldri har sett, fordi ho fangar den underliggjande strukturen. Plansjane til Jones gjer ikkje dette. Dei er ei samling former som alt finst, eit herbarium av ornament frå historia, og proposisjonane er smaksreglar destillerte ut av kva Jones og samtida hans fann harmonisk. Dei seier ikkje kvifor det mauriske ornamentet har den strukturen det har, kva krefter og teknikk og bruk og tru som forma det slik. Dei observerer at det er vakkert og foreskriv at ein gjer likeins. Skilnaden mellom å katalogisere formene og å forklare dei er heile boka sitt poeng i miniatyr. Den er ikkje at design skulle ha vorte ein generativ vitskap og feila; det er at Jones lova ei grammatikk, eit slag kunnskap, og leverte eit anna, og at faget tok det eine for det andre og har levd på forvekslinga sidan.

Kor lite ein katalog held når han blir sett på prøve, syner seg best der Jones sjølv stod midt i industrimonumentet og fekk bruke teorien sin i full skala. Då skallet til krystallpalasset stod ferdig vinteren og våren 1851, måtte nokon avgjere kva farge dei kilometerlange rekkjene av støypejarnsøyler og sprosser skulle ha. Oppgåva fall på Jones. I staden for ein verdig einsfarga tone la han ut ein systematisk fargeplan etter eit prinsipp henta frå mauriske og antikke interiør: dei tre primærfargane raudt, gult og blått, fordelte etter strenge reglar, slik at blått trekte seg attende i skuggen, gult fanga ljoset på dei framskotne kantane, og raudt låg på undersidene. Tanken var optisk: på avstand skulle dei tre smelte saman i auget til ei nøytral gråtone, og nær til gje konstruksjonen liv og djupn. Det var ikkje pynt; det var ein teori om farge sett i verk i full skala.

Striden braut laus med ein gong. Fagfolk, kritikarar og publikum delte seg; mange fann det grelt, barnsleg, framandt, andre forsvarte det; det vart halde føredrag og skrive i avisene om kven som hadde rett om fargen i taket. Jones hadde eit prinsipp og kunne grunngje valet sitt med stor lærdom. Men då striden kom, fanst det ingen grunn som batt den som var usamd. Den som fann det grelt, kunne ikkje overtydast av teorien, fordi teorien til sjuande og sist sa: slik gjorde maurarane det, og slik finn eg det harmonisk. Det er ein observasjon om kva som har vore gjort, kledd som ei lov om kva som må gjerast. Jones kunne peike på Alhambra; han kunne ikkje vise at raudt måtte ligge på undersidene på nokon måte som tvinga ein motstandar. Adrian Forty har synt korleis denne typen formstrid i reformtida like mykje var ein strid om kven som hadde rett til å avgjere smaken som om noko i sjølve fargen \autocite{forty1986}.

Det er freistande å seie at striden ikkje let seg avgjere fordi fargen mangla ei måling, slik eit rekneproblem har ho. Men det er feil veg å sjå det. Fråvêret av ei avgjerande måling gjorde ikkje fargestriden til eit mislukka rekneproblem. Det gjorde han til ein annan slags strid, ein strid om verdiar, som korkje har eller treng eit rekneproblems løysing. Striden vart likevel avgjord, og måten han vart avgjord på er heile poenget: Jones hadde oppdraget og autoriteten, og bygget skulle opne 1. mai, så fargen måtte vere på. Då dørene opna og dei seks millionane strøymde inn, gjekk dei under eit tak måla etter eit prinsipp som hadde vunne ein strid det ikkje kunne grunngje. Slik vann reforma kvar gong: ikkje med ein grunn, men med ein frist og ei stilling. Det er den same målinga skammekammeret gav, berre teken på reformatoren sjølv.

\section{Summerly: prinsippet støypt i leire}

Før Cole vart dommaren som hengde ut andre folk sine varer, var han sjølv produsent, og det avslører noko reforma helst ville gløyme: at ho fyrst prøvde å lage den gode forma, ikkje berre dømme henne. I 1847 grunnla Cole, under pseudonymet «Felix Summerly», eit lite firma han kalla Summerly's Art Manufactures. Tanken var enkel og høg: han ville få anerkjende kunstnarar til å levere teikningar til vanlege bruksting, og dermed føre kunsten og industrien saman att, lege skiljet ved å sende biletmakaren attende inn i fabrikken. Det mest kjende stykket var eit teservise. Ein tekanne teikna under Summerly-namnet vann sølvmedalje frå Society of Arts i 1846 og vart sett i produksjon hjå Mintons i Staffordshire, den same keramikkbygda Josiah Wedgwood hadde industrialisert med Etruria eit par mannsaldrar før. Cole ville at forma skulle vere «passande til føremålet»: at tuten skulle skjenke reint, at handtaket skulle høve handa, at ornamentet skulle vekse ut av bruken og ikkje liggje på som eit lass. Han prøvde å demonstrere i porselen det reforma seinare berre kunne forkynne i ord, at det finst ei rett form, og at ho kjem av at tingen skal brukast. Sparke reknar Summerly's Art Manufactures mellom dei tidlegaste medvitne forsøka på å føre kunstnarleg dømekraft inn i fabrikken \autocite{sparke2004}.

Men sjå kva som faktisk hende, for heile reforma sitt problem ligg i den eine tekannen. Dette var skiljet sett i drift på nytt, ikkje lækt. Kunstnaren teikna; arbeidarane hjå Mintons laga; Cole bar namnet og medaljen. Den «kunstnarlege» forma vart endå ei teikning levert ovanfrå til hender som ikkje teikna henne, same delinga Etruria hadde innført, no med betre smak påklistra, og det same skiljet sett i drift i den same keramikkbygda ein gong til. Og kva var grunnen til at forma var rett? Han skjenkjer reint, ja, men skjenkjer han reinare enn ein annan kanne på ein måte ein kan måle, eller skjenkjer han berre slik Cole og dei utvalde kunstnarane hans fann rett? «Passande til føremålet» høyrest ut som eit prinsipp, men spør ein kven som avgjer kva som er passande, endar ein hjå smaken til den lærde mannen att. Forty les Summerly's Art Manufactures som eit prosjekt som ville sameine kunst og industri, men som i praksis berre la eit nytt lag av borgarleg smak oppå den same gamle delinga \autocite{forty1986}. Cole prøvde å lage fundamentet i staden for å forkynne det, og enda med ein vakker tekanne som var rett fordi den rette mannen laga han. Reforma sitt mest gjennomtenkte produkt, det stykket som skulle prove at den rette forma fanst, kviler til sjuande og sist på kven som heldt penselen.

\section{South Kensington: katalogen gjord til bygning}

Reforma gav frå seg éin ting som varte lenger enn alle dommane og alle katalogane: ein institusjonell modell for korleis ein skulle lære den industrielle forma. Etter at Department of Practical Art vart oppretta i 1852, voks det kring samlingane som flytta til South Kensington fram eit heilt apparat, eit museum som seinare vart Victoria and Albert, ein sentralskule, og eit nett av provinsskular bundne til hovudstaden gjennom felles pensum, felles eksamenar og felles autoriserte førebilete. Dette er «South Kensington-systemet», og det vart eksportert; mykje av den måten designskular framleis er organiserte på, kan førast attende hit.

Sjølve museet var den geniale og den feilslegne ideen på éin gong. Cole-krinsen samla inn gjenstandar, historiske og samtidige, gode døme og skrekkdøme, og stilte dei ut ikkje for å glede, men for å lære: museet var ein lærebok i tre dimensjonar, ei samling autoriserte former studenten skulle gå mellom, teikne av og innprente seg. Tanken var at god form kunne overførast ved synet av rette døme, slik eit språk kunne lærast ved å lese rette setningar. Men eit museum av døme er ein katalog gjord til bygning, same forma for kunnskap som grammatikken til Jones, no med veggar og opningstider. Det viser deg kva som er rett; det seier deg ikkje kvifor. Studenten som vandra gjennom South Kensington og teikna av den rette vasen, lærte å kjenne att og gjengje ein kanon. Han lærte ikkje ein teori som kunne fortelje han kvifor kanonen var rett, eller late han avgjere ei ny form museet aldri hadde sett. Det er det same skulane underviste i smått: kopiering og innprenting, lauget sin metode flytta frå verkstaden til klasserommet, ein dom utan ein grunn under dommen.

Forty peikar på spenninga systemet aldri løyste: det skulle danne smaken til ein heil nasjon, men kunne ikkje grunngje smaken på noko som batt den som var usamd, og vart difor til sjuande og sist ei øving i å innprente autoritetens preferansar i stor skala \autocite{forty1986}. John Heskett strekar under at perioden bygde imponerande institusjonar for design lenge før det fanst noko fagleg innhald å fylle dei med, at huset for utdanninga reiste seg før kunnskapen som skulle undervisast i det \autocite{heskett1980}. Det er den same baklengs rekkjefølgja som i industrien, no flytta inn i skulestova: fyrst museet, eksamenen, pensumet, og inni dei eit tomrom der fundamentet skulle ha vore. Og fordi modellen var så lett å kopiere og så grei å eksportere, bar han tomrommet sitt vidare til kvar skule som tok han etter.

\section{Skallet ingen kunne lese}

Det ligg ein ironi i hjartet av 1851 som reformatorane aldri heilt fekk auge på. Då komiteen for utstillinga i 1850 hadde forkasta hundretals innsende forslag, tunge murbygg som ville teke år å reise og aldri vorte ferdige i tide, kom det inn eit utkast frå ein mann som ikkje var arkitekt i det heile: Joseph Paxton, gartnaren ved Chatsworth, som hadde bygd drivhus for hertugen av Devonshire. Han teikna utkastet, etter eiga forteljing, på eit trekkpapir under eit kjedeleg jernbanestyremøte i juni 1850. Det han la fram, var ikkje eit hus, men eit drivhus blåst opp til katedralstorleik, eit skall av støypejarn og glas, sett saman av standardiserte delar i ein modul på tjuefire fot, som kunne masseproduserast, fraktast og skruast saman på staden.

Tala held seg svimlande. Bygget reiste seg på under tjuetre veker; det dekte om lag 92 000 kvadratmeter; det slukte i storleiksorden 300 000 glasplater; og det var høgt nok til at dei store almetrea i parken kunne stå urørte inni, under taket. Då det opna 1. mai 1851, strøymde det inn folk i eit tal verda ikkje hadde sett: før det stengde i oktober hadde kring seks millionar vitjingar passert, opp mot hundre tusen på dei travlaste dagane. Satirebladet \textit{Punch} gav det namnet som feste seg, krystallpalasset, meint som ein spøk, motteke som ei sanning. Og dette bygget, som folk straks las som det vakraste og mest framtidsvende dei hadde sett, braut på nesten kvart punkt med prinsippa Cole-krinsen forkynte inni det. Det skjulte ikkje konstruksjonen under pynt; det var konstruksjon, naken og repetert. Det etterlikna ikkje eitt materiale i eit anna; det viste jarnet som jarn og glaset som glas. Det hadde inga unødig ornamentering, fordi det korkje hadde tid eller råd til noko slikt. Det gjorde alt Pugin og Jones bad om, og det gjorde det utan å meine det.

Her er kjernen, og han er lett å mistyde. Poenget er ikkje at ingeniørkulturen og hagebruket var ein høgare disiplin design skulle ha etterlikna, og heller ikkje at Paxton gjorde betre design enn reformatorane. Paxton dreiv ikkje med smak i det heile. Poenget er at forma hans hadde ein grunn utanfor den som dømde. Glaset og jarnet, modulen, fristen og prisen var krefter som verka på forma same kven som heldt blyanten, og dei batt forma slik ingen autoritet i Cole-krinsen kunne binde henne. Reformatorane leitte etter forma sitt fundament i smaken og moralen, medan bygget kring dei viste at forma sin verkelege grunn låg i krefter: i kva glaset og jarnet ville og ikkje ville, i kva modulen kravde, i kva tida tillét. Dei stod bokstaveleg inni eit svar dei ikkje kunne lese. Dei dømde tinga på borda og overså bygget over hovuda sine.

Éin mann las det motsett, og det er lærerikt at han stod i same bygg som Cole-krinsen og likevel vart aldri følgd. Gottfried Semper, tysk arkitekt i London-eksil frå 1850 til 1855, gjekk gjennom den same utstillinga, men i staden for å dømme smaken spurde han kvar forma kjem frå, og svaret hans peika nedover, mot materialet og teknikken. I \textit{Der Stil} la han fram tanken om at dei grunnleggjande formene i byggjekunsten spring ut av nokre urhandverk, vevnaden, keramikken, tømringa, muringa, og at meininga i forma heng saman med den tekniske handlinga ho voks ut av \autocite{semper1860stil}. Her var eit forsøk på å forklare forma ut frå krefter som verka på henne, å føre henne attende til ein grunn utanfor smaken til den som dømer. Det er det næraste hundreåret kom eit fundament, og det vart aldri det. Semper sitt prosjekt vart lese og dels misforstått, men det vart éi stemme mellom mange, ein teori blant teoriar, integrert i ingen felles disiplin. Heskett peikar på kor symptomatisk det er at perioden produserte glimrande einskildtenkjarar, men ingen kumulativ teori dei kunne byggje saman på \autocite{heskett1980}. Fundamentet, om det fanst nokon stad i Hyde Park i 1851, stod i jarnsøylene, ikkje i katalogane, og ingen i hallane kunne lese søylene.

\begin{kasus}{Raudt, gult, blått: fargestriden i Hyde Park, våren 1851}
Owen Jones står oppe i ei stillasrigg under glastaket, med malarane sine under seg og kilometer av jarn framfor seg, og dirigerer kvar kulør på kvar flate etter ein plan han har i hovudet frå Alhambra. Blått i skuggen, gult på ljoskanten, raudt på undersida. Under han samlar kritikarane seg og peikar: dette er for grelt, dette er ikkje verdig eit nasjonalt monument. Jones svarar med teorien sin. Båe partar har rett innanfor sitt eige rom, og ingen kan tvinge den andre, fordi det ikkje finst nokon grunn over dei begge. Striden vart likevel avgjord, ikkje av ein teori, men av at Jones hadde oppdraget og at bygget skulle opne 1. mai. Då dei seks millionane strøymde inn, gjekk dei under eit tak måla etter eit prinsipp som hadde vunne ein strid det ikkje kunne grunngje. Striden trong aldri eit rekneproblems løysing; han var ein verdistrid. Men eit fag som avgjer verdistridane sine med frist og stilling, har innrømt at det ikkje kan grunngje forma si.
\end{kasus}

\vignett

Reformatorane sette seg føre å gje forma eit fundament og kom heim med fire ting som alle ser ut som ein grunn og ingen er det: ein dom, ein moral, ein katalog og ein institusjon. Kvar gong dommen vart sett på prøve, av eit publikum som likte sine eigne ting, av produsentar som protesterte, av ein motstandar i fargestriden, hadde han ingenting å falle tilbake på utan meir autoritet. Den vakraste forma i Hyde Park stod i søylene over hovuda deira og vart aldri lesen, fordi ho var teikna av tvang og dei leitte etter ein smak. Men reforma var ikkje over. Ho skulle straks finne ein meir lidenskapleg skapnad, hjå menn som ikkje ville reformere den industrielle forma, men avvise heile maskina som hadde skapt henne, og kalle handverket tilbake som botemiddel. Det er den fyrste store rørsla, og ho byrjar med ein sint mann og eit kapittel om gotikkens natur.

\laeringsmal{\begin{itemize}
\item Plassere den store utstillinga i 1851 som det fyrste augneblinket nokon stod andsynes den industrielle forma i full breidd og spurde etter ein standard for henne.
\item Skilje mellom ein dom, ein moral, ein katalog og eit fundament, og bruke skiljet på Cole, Pugin og Jones.
\item Kjenne att sviktsignaturen: ein strid som berre kan avgjerast med rang, frist eller stemme, og vise korleis han verkar i skammekammeret og i fargestriden.
\item Gjere greie for kvifor ironien i skallet ikkje gjer ingeniørfaget til ein mal design feila ved, men syner at Paxton si form hadde ein grunn utanfor den som dømde.
\end{itemize}}

\ovingar{\begin{enumerate}
\item Owen Jones kalla verket sitt ein grammatikk. Ein grammatikk er generativ; ein katalog er det ikkje. Vel ein samtidig «designkanon» eller stilguide og avgjer om han er grammatikk eller katalog. Korleis ser ein skilnaden?
\item Pugin grunngav forma med teologi, Loos med ein rangstige av folkeslag. Begge kledde ein smak som ei lov. Finn eit moderne døme der ein estetisk preferanse blir presentert som ein naudsynt konsekvens av noko utanfor smaken.
\item Sett opp ditt eige «skammekammer» med tre samtidige gjenstandar du finn dårlege. Skriv dommen og grunngjevinga. Be ein medstudent som er usamd avvise grunngjevinga. Kva held, og kva fell?
\item Fargestriden var ein verdistrid utan eit rekneproblems løysing. Diskuter: finst det formstridar som faktisk \emph{har} ei avgjerande måling, eller endar dei alle hjå verdiar til slutt?
\end{enumerate}}

\vidarelesing{\fullcite{auerbach1999}; \fullcite{jones1856grammar}; \fullcite{pugin1841true}; \fullcite{pugin1836contrasts}; \fullcite{semper1860stil}; \fullcite{loos1908ornament}; \fullcite{sparke2004}; \fullcite{forty1986}.}
